\documentclass[a4paper,12pt]{report}
\usepackage{mathptmx}                   % Usa Times (parecida a Times New Roman)
\usepackage{setspace}                   % Para interlineado
\onehalfspacing                         % Establece interlineado 1.5
\usepackage{xcolor}                     % Para colores personalizados
\usepackage[table]{xcolor}              % Para los colores en las tablas
% \usepackage[table,xcdraw]{xcolor}     % Para incluir opciones de color en la tabla
\usepackage{graphicx}                   % Para incluir imágenes
\usepackage{pdfpages}                   % Para incluir pdf
\usepackage{geometry}                   % Para ajustar márgenes
\usepackage{titlesec}                   % Para personalizar títulos
\usepackage{fancyhdr}                   % Para encabezados y pies de página
\usepackage{hyperref}                   % Para incluir URLs
\usepackage{csquotes}                   % Para las citas inline 
\usepackage[backend=bibtex]{biblatex}   % Para añadir bibliografía con BibLaTeX
\usepackage{float}                      % Para las tablas

% Configuración de márgenes
\geometry{left=3cm, right=3cm, top=2.5cm, bottom=2.5cm}
\graphicspath{{resources/}{../pdf/}}

% Personalización de títulos
\titleformat{\chapter}[display]
  {\Huge\bfseries}{}{0pt}{\Huge}

% Encabezados y pies de página
\pagestyle{fancy}
\fancyhead{}
\fancyfoot{}
\fancyfoot[C]{\thepage}

% Partición de palabras en textos justificados
\hyphenation{je-fa-tu-ra}
\hyphenation{LABORA}
\hyphenation{ma-nu-fac-tu-ra}
\hyphenation{pro-fe-sio-nal}
\hyphenation{pro-fe-sio-na-les}
\hyphenation{co-rres-pon-dien-tes}
\hyphenation{pro-xi-mi-da-des}

% Indicar ficheros con la bibliografía
\addbibresource{biblio.bib}

% Configuración de hipervínculos
\hypersetup{
    colorlinks=true,
    linkcolor=black,    % Enlaces internos (a secciones, figuras, etc.)
    urlcolor=gray,      % Enlaces externos (URLs)
    citecolor=black,    % Referencias/citas
    pdftitle={Memoria Prácticum II},
    pdfauthor={Nombre Apellido},
    pdfsubject={Prácticas en centro educativo}
}

\begin{document}
    \begin{titlepage}
    \begin{center}
        \vspace*{2cm}
        
        \includegraphics[width=0.5\textwidth]{resources/logo_ua.png}\\[1cm]
        
        \textcolor{gray}{\Huge \textbf{Memoria del Prácticum II}}\\[1cm]
        \rule{10cm}{0.4mm} \\[1cm]
        
        {\Large Máster en Profesorado de Educación Secundaria Obligatoria y Bachillerato, Formación Profesional y Enseñanza de Idiomas}\\[1cm]
        
        \textbf{Especialidad:} Especialidad 3. Formación Profesional Cuerpo Técnico\\[0.5cm]
        \textbf{Alumno/a:} Lorena González Núñez de Arenas\\[0.5cm]
        \textbf{Curso académico:} 2024-2025\\[0.5cm]
        \textbf{Centro de Prácticas:} CIPFP Canastell\\[0.5cm]
        \textbf{Tutor/a del centro:} Antonio Verdejo Borja\\[0.5cm]
        \textbf{Tutor/a de la Facultad:} Raúl Gutierrez Fresneda\\[0.5cm]
        
        \vfill
        \textcolor{gray}{\Large \textbf{Universidad de Alicante}}\\
        \textcolor{gray}{Facultad de Educación}\\
        \textcolor{gray}{San Vicente del Raspeig, Alicante}\\
        \textcolor{gray}{Fecha: \today}
    \end{center}
\end{titlepage}
    % Declaración de Responsabilidad
\newpage
\section*{DECLARACIÓN DE RESPONSABILIDAD Y AUTORÍA}
D/Dª.: Lorena González Núñez de Arenas, con DNI 48620705Q, 
estudiante del Máster universitario en profesorado de educación secundaria obligatoria y bachillerato, formación profesional y enseñanza de idiomas, 
de la Universidad de Alicante, realizado en el período del 20 de enero de 2025 al 20 de febrero de 2025. \\


DECLARA QUE:\\
La Memoria del Prácticum denominado\\
Memoria del Prácticum I\\
ha sido desarrollado respetando los derechos intelectuales de terceros, conforme las citas que 
constan en las páginas correspondientes y cuyas fuentes se incorporan en la bibliografía, así 
como cualquier otro derecho, por ejemplo de imagen que pudiese estar sujeto a protección 
del copyright.\\

En virtud de esta declaración, afirmo que este trabajo es inédito y de mi autoría, por lo que me 
responsabilizo del contenido, veracidad y alcance de la Memoria del Prácticum, y asumo las consecuencias administrativas y jurídicas 
que se deriven en caso de incumplimiento de esta declaración.\\

Para que así conste, firmo la presente declaración en \\
Alicante, a \today. \\

Fdo.:  \
\vspace{10\baselineskip}\

Este documento formará parte de la memoria de los Practicum o TFG o TFM correspondiente y 
será la primera página de los mismos.\\

\textsuperscript{*}	\textit{Documento aprobado en Junta de Facultad el 19 de octubre de 2017.}\\


    \renewcommand{\contentsname}{Índice}
    \tableofcontents
    \newpage

    % RESUMEN
\chapter{Resumen}

Esta memoria recoge la experiencia y las intervenciones realizadas durante el Prácticum II en el CIPFP Canastell, dentro del módulo de \textit{Instalaciones Domóticas} del segundo curso del Grado Medio en Instalaciones Eléctricas y Automáticas. La intervención se centró en una situación de aprendizaje que abordó los riesgos de seguridad en la domótica de bajo coste, permitiendo al alumnado analizar vulnerabilidades en dispositivos IoT y aplicar buenas prácticas de ciberseguridad.

A lo largo del período de prácticas, mi rol principal fue de observadora, ya que mi formación en Ingeniería Informática no estaba alineada con el contenido eléctrico del ciclo. Sin embargo, esto no impidió que pudiera desarrollar una intervención concreta en un área más afín a mi perfil: la seguridad en domótica de bajo coste. 

La situación de sprendizaje titulada "Domótica de Bajo Coste: ¿Comodidad Inteligente o Riesgo Silencioso?", se llevó a cabo el 6 de febrero de 2025. En esta sesión, el alumnado analizó vulnerabilidades de dispositivos IoT de bajo coste y aplicó buenas prácticas en ciberseguridad. La actividad combinó una parte teórica donde se explicaron los riesgos y ataques más comunes en redes domóticas, y una parte práctica en aula-taller, en la que los estudiantes experimentaron con la configuración y seguridad de un \textit{enchufe inteligente Shelly Plug} y un \textit{router WiFi}, identificando riesgos potenciales y simulando ataques de fuerza bruta.

Se utilizaron estrategias didácticas activas, fomentando la participación del alumnado mediante dinámicas interactivas y herramientas digitales. La evaluación se realizó a través de una prueba tipo test y la observación del desempeño práctico.

Los resultados evidencian que la actividad resultó altamente efectiva en la concienciación sobre la ciberseguridad en el ámbito domótico, con una participación notable en la parte práctica. No obstante, se identificaron dificultades en la motivación de ciertos alumnos y en la gestión del tiempo en la actividad práctica. 

En conclusión, esta experiencia ha permitido desarrollar competencias pedagógicas clave, mejorando la capacidad de planificación, adaptación metodológica y gestión del aula en un entorno de Formación Profesional.

% \textbf{Palabras clave}: Prácticum II, Formación Profesional, domótica, ciberseguridad, enseñanza práctica, aprendizaje basado en problemas.



% Durante la observación en el aula, pude identificar dinámicas interesantes en la interacción profesor-alumno, así como detectar dificultades individuales. Entre los casos destacados:
% \begin{itemize}
%     \item Un alumno de \textit{1º GM} mostraba signos de falta de concentración y comentó que no podía dormir bien. Tras una breve conversación en la que le sugerí buscar ayuda profesional, retomó la práctica con mayor enfoque.
%     \item Otro estudiante de \textit{1º GM} presentó una descripción ambigua en una práctica. Tras hacerle reflexionar sobre la claridad de su explicación desde el punto de vista de un cliente, modificó su respuesta.
%     \item Un alumno me llamó \textit{"profe"} y me preguntó sobre disyuntores. Tras explicarle que no podía ayudarle porque no era mi área, entablamos una conversación en la que cuestionó por qué me habían asignado a un módulo eléctrico, reflejando la falta de alineación entre mi perfil y la asignación de prácticas.
% \end{itemize}

% La evaluación de la actividad evidenció que la parte práctica generó gran interés en el alumnado, destacando la importancia de incluir más contenido de ciberseguridad en la formación de Instalaciones Domóticas. A pesar de las limitaciones derivadas de mi perfil técnico, esta experiencia me permitió reflexionar sobre el papel del docente en FP, el enfoque metodológico en enseñanza práctica y la importancia de la motivación y adaptación de los contenidos a las necesidades del alumnado.


    % INTRODUCCIÓN
\chapter{Introducción}
El presente documento recoge la experiencia desarrollada durante el período del Prácticum II correspondiente al \textit{Máster en Profesorado de Educación Secundaria Obligatoria y Bachillerato, Formación Profesional y Enseñanza de Idiomas}, llevado a cabo en el CIPFP Canastell, centro integrado de Formación Profesional ubicado en San Vicente del Raspeig. Este centro destaca por su amplia oferta formativa adaptada al contexto industrial y empresarial local, con especial énfasis en las familias profesionales relacionadas con la electricidad y electrónica, entre otras.

Durante mi estancia en el centro, mi rol principal fue el de observadora, participando de manera puntual en algunas intervenciones pedagógicas. Este enfoque fue consecuencia directa de la falta de adecuación entre mi perfil formativo, Ingeniería Informática, y el contenido curricular mayoritariamente eléctrico de los ciclos en los que se desarrolló la práctica docente. Esta situación no solo condicionó mi nivel de participación directa en la docencia, sino que también generó una reflexión crítica sobre la importancia de alinear los perfiles de formación del profesorado en prácticas con las asignaturas asignadas.

Mi intervención principal tuvo lugar en el módulo de \textit{Instalaciones Domóticas}, impartido en el segundo curso del Grado Medio en Instalaciones Eléctricas y Automáticas, durante una sesión realizada el 6 de febrero de 2025. Esta actividad, titulada ``Domótica de Bajo Coste: ¿Comodidad Inteligente o Riesgo Silencioso?'', permitió al alumnado explorar aspectos fundamentales de ciberseguridad aplicados a la domótica de bajo coste, tema más próximo a mi ámbito de formación.

En dicha unidad de trabajo se desarrollaron diversas dinámicas que facilitaron la interacción y participación activa del alumnado. Entre éstas destacó la utilización de herramientas digitales como un código QR que les permitió contribuir a una nube de palabras sobre riesgos de seguridad en domótica. Este tipo de recursos sirvió como rompehielo para generar un ambiente propicio para el aprendizaje activo. Además, la práctica consistió en configurar un enchufe inteligente \textit{Shelly Plug}, conectarlo a una red WiFi con seguridad débil y simular ataques informáticos básicos, actividades que generaron un gran interés en las/os estudiantes, evidenciado en la alta participación y la profundidad de las reflexiones posteriores.
Cabe mencionar que hubo un estudiante que fue capaz de cambiarle el nombre a la red, haciendo que no nos pudiéramos conectar al enchufe inteligente. Al finalizar la clase, se acercó para preguntarme si había ``molestado'' con su acción, a lo que le respondí que no, que había sido una acción muy interesante y que era una forma de aprender.

Por otro lado, la observación directa de dinámicas de enseñanza en el aula de 1º GM, específicamente en el módulo de \textit{Automatismos Industriales}, aportó experiencias valiosas sobre la realidad cotidiana del profesorado en FP. Durante estas sesiones observé situaciones que enriquecieron mi visión acerca de la complejidad de la labor docente, especialmente cuando se trata de estudiantes con diferentes necesidades y motivaciones. Entre estas situaciones destaco las siguientes anécdotas registradas en mi cuaderno de bitácora:

\begin{itemize} %[leftmargin=*, label=\textbullet]
    \item En una sesión práctica, observé a un alumno de 1º GM aparentemente desconectado de la actividad. Al acercarme, me comentó que tenía problemas de sueño y dificultades para concentrarse. Aproveché la oportunidad para recomendarle que buscara ayuda profesional, sugiriéndole acudir al médico o a un especialista. Tras esta breve conversación y sin mediar palabras, noté un cambio positivo en su actitud, retomando la actividad con mayor atención y disposición.
    
    \item Otra situación se produjo cuando un alumno presentó una descripción escueta y ambigua en la explicación de un automatismo. Al pedirme el docente que la revisara, le hice ver al estudiante que su descripción no era suficiente desde la perspectiva de un cliente, destacando la importancia de una buena comunicación técnica. Aunque inicialmente hubo cierta resistencia, finalmente aceptó mejorar su respuesta gracias a la insistencia del profesor titular.
    
    \item También recuerdo con especial ilusión cuando uno de los alumnos se refirió a mí como ``profe'' para preguntarme sobre una cuestión técnica relacionada con disyuntores. Dado que era un tema fuera de mi área de especialización, le aclaré que no podía ayudarle y le remití al profesorado titular. Esta situación derivó en una conversación en la que el alumno, con notable lógica y observación, cuestionó por qué me habían asignado un módulo de electricidad siendo ingeniera informática, mostrando la clara falta de adecuación en la asignación de mis prácticas.
    
    \item En otras sesiones, observé situaciones diversas, como alumnas/os que mostraban estrés por exámenes o tareas, y pude apreciar cómo el profesorado abordaba estas situaciones desde una perspectiva no solo académica, sino también emocional, enfatizando la importancia de un acompañamiento integral a las/os estudiantes en FP.
\end{itemize}

Estas anécdotas reflejan no solo aspectos concretos de la realidad cotidiana del aula, sino también la importancia de la figura docente en Formación Profesional como agente clave en el desarrollo integral del alumnado, desde lo académico hasta lo personal y emocional.

Finalmente, cabe destacar que esta experiencia ha supuesto una oportunidad única para reflexionar sobre las dinámicas didácticas empleadas en FP, la necesidad de adaptar los contenidos a las características y motivaciones del alumnado, así como para entender la trascendencia de ofrecer un contexto de aprendizaje seguro y propicio para la participación. A pesar de las limitaciones derivadas de mi formación inicial, esta experiencia ha sido enriquecedora en múltiples dimensiones, contribuyendo notablemente a mi formación pedagógica y proporcionando aprendizajes que serán aplicables en futuros contextos docentes.

    % CONTEXTUALIZACIÓN
\chapter{Descripción y valoración de la situación educativa concreta}

Durante el período de prácticas del Prácticum II, mi actividad principal se ha desarrollado en el \textit{CIPFP Canastell}, un centro integrado de Formación Profesional ubicado en San Vicente del Raspeig. En este contexto, he participado en tres módulos de formación pertenecientes a dos ciclos distintos del ámbito de la electricidad y la mecatrónica (Instalación y Mantenimiento):

\begin{itemize}
    \item \textbf{Ciclo de Grado Medio de Instalaciones Eléctricas y Automáticas}:
    \begin{itemize}
      \item Módulo ``Automatismos Industriales'' (1º Curso)
      \item Módulo ``Instalaciones Domóticas'' (2º Curso)
    \end{itemize}
    \item \textbf{Ciclo de Grado Superior de Mecatrónica Industrial}:
    \begin{itemize}
      \item Módulo ``Sistemas Eléctricos y Electrónicos'' (1º Curso)
    \end{itemize}
  \end{itemize}\

La participación se centró principalmente en el módulo de \textit{Instalaciones Domóticas} de segundo curso del Ciclo Formativo de Grado Medio de Instalaciones Eléctricas y Automáticas, y de forma observacional en los otros dos módulos.


\section{Características del alumnado}

Durante el período de prácticas, correspondiente al segundo trimestre del curso académico 2024-25, he podido observar una diferencia significativa entre el número de estudiantes matriculados oficialmente y el número real de asistentes a clase en los diferentes módulos. Según me explicó el tutor, esta situación es habitual en los ciclos formativos, especialmente en primero del Grado Medio, donde el primer trimestre actúa como una especie de filtro natural. Durante ese periodo inicial, suelen concentrarse los alumnos más conflictivos o con menor compromiso académico, lo que convierte las clases en un entorno más complejo para la docencia. Sin embargo, a partir del segundo trimestre, muchos de esos estudiantes se van descolgando progresivamente, lo que conlleva una mejora notable del clima en el aula y una mayor estabilidad en los grupos. Lo cual no siginifica que no haya alumnado conflictivo, sino que la mayoría de ellos ya no asisten a clase.

Mi experiencia como docente en prácticas, al desarrollarse durante este segundo trimestre, se ha dado en un contexto ya más consolidado, con grupos más reducidos, cohesionados y participativos. Esto ha influido positivamente en el desarrollo de las sesiones, especialmente en aquellas en las que he tenido una intervención directa.

En el módulo de \textit{Instalaciones Domóticas} (2º curso del Ciclo Formativo de Grado Medio de Instalaciones Eléctricas y Automáticas), asistían regularmente alrededor de quince estudiantes, la mayoría, mayores de edad, a partir de los diecisiete años. El grupo mostraba una composición heterogénea, tanto en lo académico como en lo experiencial: algunos alumnos accedieron tras finalizar la ESO, otros por prueba de acceso y varios con experiencia previa en el ámbito técnico o laboral. Aunque el interés general por los contenidos era moderado, existían claras diferencias en cuanto a implicación, motivación y ritmo de aprendizaje.

En el módulo de \textit{Automatismos Industriales} (1º curso del mismo ciclo), que seguí desde un enfoque observacional, la asistencia se situaba entre los diecisiete y diecinueve estudiantes, con edades a partir de los dieciséis años. Este grupo se encontraba todavía en proceso de adaptación al entorno y las exigencias de la Formación Profesional, presentando dificultades organizativas, menor autonomía y dispersión en las sesiones prácticas. A pesar de la estabilización tras el primer trimestre, el profesorado destacaba la necesidad de un acompañamiento más intensivo. En este grupo, además, se identificaban casos de alumnado con necesidades educativas específicas, como el de una alumna procedente de la ESO que evidenciaba tanto dificultades académicas como personales, y que requería una atención empática, aunque no habían estrategias de inclusión por parte del equipo docente de manera formal, ya que las adaptaciones en Formación Profesional son de acceso. Sin embargo, sí que había predisposición del tutor y el otro profesor con el que hacía desdoble, de implicarse y realizar apoyo emocional a esta alumna.

Finalmente, en el módulo de \textit{Sistemas Eléctricos y Electrónicos} (1º curso del Ciclo Formativo de Grado Superior de Mecatrónica Industrial), también desde un rol de observadora, la asistencia era más reducida, con unos doce estudiantes a partir de los dieciocho años. Al tratarse de un ciclo superior, el perfil del alumnado era más técnico y profesionalizado. Sin embargo, no estaban exentos de dificultades, especialmente aquellos que compatibilizaban los estudios con un empleo. Esta doble carga provocaba tensiones a la hora de programar exámenes u otras actividades, generando la necesidad de una mayor flexibilidad institucional para responder a las circunstancias personales del alumnado.

En conjunto, el alumnado con el que tuve contacto reflejaba la diversidad y complejidad propias de los ciclos de Formación Profesional, tanto en lo académico como en lo personal, y me permitió observar y reflexionar sobre diferentes estrategias de gestión del aula, acompañamiento individualizado e inclusión educativa.

A continuación se presentan algunos casos concretos que ilustran la diversidad del alumnado y las dinámicas observadas en el aula:

\textbf{Caso 1: Alumno con problemas de sueño}  
Durante una práctica, observé que un alumno estaba distraído, manipulando un tornillo sin interactuar con la actividad. Al preguntarle, mencionó que no podía pensar ni dormir. Le sugerí acudir a un especialista si tenía dificultades de descanso. Tras la conversación, retomó la actividad.

\textbf{Caso 2: Explicación insuficiente en una práctica}  
Un alumno presentó una respuesta muy escueta en la descripción de un automatismo. Comenté que, desde el punto de vista del cliente, no quedaba claro su funcionamiento. Aunque inicialmente protestó, modificó su respuesta.

\textbf{Caso 3: Reflexión sobre la asignación de prácticas}  
Un alumno me preguntó sobre un problema con disyuntores. Le expliqué que no podía ayudarle y que preguntara a los profesores. Luego, me preguntó por qué, siendo informática, estaba en un ciclo de electricidad. Esto reflejó la falta de coherencia en la asignación de prácticas.



\section{Condiciones físicas del aula}

Las sesiones en las que he participado, tanto en calidad de observadora como en mi intervención activa, se han desarrollado en diferentes espacios del CIPFP Canastell, adaptados a las necesidades técnicas y metodológicas de cada módulo. Todos ellos se corresponden con aulas-taller, con una distribución funcional del espacio que facilita el aprendizaje práctico, aunque también presentan ciertas limitaciones derivadas de la antigüedad de las instalaciones.

El módulo de \textit{Instalaciones Domóticas} se imparte en un aula-taller especializada, en cuyo almacén se conservan diversos materiales destinados a las prácticas del alumnado, como componentes domóticos (enchufes inteligentes, bombillas, routers, etc.) y otros elementos como dispositivos KNX. Para el desarrollo de la situación de aprendizaje centrada en la ciberseguridad en redes domóticas, configuré un router independiente, sin conexión a Internet, que generaba un SSID denominado \textit{RedDomotica}. Este entorno simulado permitió al alumnado experimentar con la conexión y configuración de dispositivos inteligentes en condiciones seguras. La disposición del aula favorece el trabajo individual y por parejas, y se complementa con el uso de una pizarra digital de última generación, modelo \textit{SYNETECH Advance Aquarius} (A7532), incorporada este mismo curso gracias a los fondos europeos destinados a la mejora de la Formación Profesional. Esta herramienta ha supuesto una mejora significativa en la dinámica de clase, ya que permite, entre otras funcionalidades, compartir las anotaciones de la pizarra mediante un código QR, lo que facilita que el alumnado pueda centrarse en la explicación sin preocuparse por copiar apuntes en tiempo real.

En el módulo de \textit{Automatismos Industriales}, las prácticas se llevan a cabo en un taller con mesas técnicas equipadas con cuadros eléctricos, pulsadores, relés y autómatas programables. Durante una de las sesiones, detectamos que algunos tableros estaban sueltos debido a la falta de tornillos, por lo que mi tutor y yo procedimos a su reparación. Mientras tanto, el alumnado continuó con las prácticas que tenían pendientes. Este tipo de incidencias materiales es relativamente frecuente en el centro, y se entienden dentro del contexto de unas instalaciones con más de 40 años de antigüedad. De hecho, algunos elementos del mobiliario, como taburetes o mesas, conservan aún el etiquetado original de fabricación de los años 80, lo que ilustra bien la necesidad de una renovación progresiva de infraestructuras.

En el módulo de \textit{Sistemas Eléctricos y Electrónicos}, las clases se desarrollaban en un aula-taller adecuada en cuanto a espacio, aunque sin demasiadas novedades en cuanto a equipamiento. El número reducido de estudiantes favorecía un ambiente tranquilo y propicio para la formación, aunque las sesiones de teoría prolongada dificultaban en ocasiones el mantenimiento de la atención. En este caso, la estructura tradicional del aula contrastaba con las mejoras tecnológicas observadas en otros espacios del centro.

En conjunto, puede afirmarse que los espacios utilizados están ajustados a las necesidades de los distintos módulos, si bien el desgaste acumulado con el paso de los años y el uso intensivo hacen necesaria una inversión sostenida para mantener su funcionalidad. Este objetivo forma parte de las prioridades estratégicas recogidas tanto en el \textit{Manual de Calidad} como en el \textit{Proyecto de Dirección} del centro.

    % SITUACIÓN DE APRENDIZAJE
\chapter{Situación de Aprendizaje: Domótica de Bajo Coste}
\section{Justificación y Contexto}
El uso de dispositivos domóticos de bajo coste está en auge, pero la seguridad sigue siendo un aspecto desatendido. Esta SA busca concienciar a los alumnos sobre las vulnerabilidades de estos sistemas.

\section{Objetivos}
\begin{itemize}
    \item Comprender los riesgos de seguridad en la domótica de bajo coste.
    \item Identificar vulnerabilidades en dispositivos IoT.
    \item Aplicar buenas prácticas de configuración segura.
\end{itemize}

\section{Contenidos}
\begin{itemize}
    \item Seguridad en redes WiFi y dispositivos inteligentes.
    \item Métodos de ataque: fuerza bruta, sniffing, spoofing.
    \item Configuración de dispositivos domóticos seguros.
\end{itemize}

\section{Desarrollo de la sesión}
\textbf{1. Teoría (1h)}  
Explicación sobre vulnerabilidades en domótica y medidas de protección.

\textbf{2. Práctica (1h 45min)}  
\begin{itemize}
    \item Configuración de un enchufe inteligente Shelly Plug.
    \item Simulación de ataques a una red WiFi insegura.
    \item Análisis de permisos de aplicaciones domóticas.
\end{itemize}

\section{Evaluación}
\begin{itemize}
    \item Participación en la sesión (30\%).
    \item Realización de la práctica (40\%).
    \item Examen tipo test (30\%).
\end{itemize}
    % ANÁLISIS Y VALORACIÓN DE LA INTERVENCIÓN
\chapter{Análisis y valoración de la intervención}

\section{Reflexión sobre la intervención como docente en prácticas}

Mi participación durante el Prácticum II ha sido fundamentalmente de observación activa, aunque también he podido intervenir de forma puntual como apoyo a la docencia y, de forma más significativa, en la planificación, desarrollo y evaluación de la unidad de trabajo sobre seguridad en dispositivos domóticos. Esta intervención me ha permitido poner en práctica mis conocimientos técnicos desde una perspectiva didáctica, reflexionar sobre el funcionamiento real de un aula de Formación Profesional y adaptar mi comunicación y estrategias a un perfil de alumnado diverso.

A lo largo del periodo de prácticas, observé dinámicas muy diferentes entre los grupos, especialmente entre primer y segundo curso del ciclo de grado medio. En algunos casos, la desmotivación o el uso inadecuado del móvil dificultaban el ritmo de las clases, mientras que en otros momentos la implicación era notable, sobre todo cuando las actividades se conectaban con su entorno cotidiano o se planteaban en formato de reto.

La experiencia más significativa fue sin duda la unidad de trabajo que diseñé y llevé a cabo en el módulo de \textit{Instalaciones Domóticas}. Su planificación, desarrollo y evaluación me permitieron explorar todo el ciclo didáctico y afrontar imprevistos reales (como problemas técnicos o intervenciones espontáneas del alumnado), lo que enriqueció enormemente mi formación como futura docente.

\section{Valoración de las propuestas pedagógicas del departamento}

El Departamento de Electricidad y Electrónica muestra un compromiso activo con la mejora continua y con la aplicación de metodologías que favorezcan el aprendizaje práctico. A pesar de que el enfoque dominante sigue siendo la clase magistral tradicional, se observan esfuerzos por incorporar metodologías activas, como el aprendizaje basado en retos (ABR), especialmente en módulos como Robótica.

Sin embargo, también se detectan limitaciones estructurales y de formación docente que condicionan la implantación de enfoques más innovadores. Por ejemplo, una profesora comentó que aplicaba ABR sin haber recibido formación específica, lo que puede limitar su potencial. Aun así, se aprecia una buena disposición del profesorado a colaborar, compartir materiales y debatir sobre estrategias didácticas, lo cual es una fortaleza clara del departamento.

\section{Valoración de la programación de aula y la situación de aprendizaje}

La programación del módulo de Instalaciones Domóticas establece un marco coherente con el currículo oficial y recoge una organización de contenidos por unidades de trabajo. Mi unidad de trabajo se diseñó de forma complementaria a la programación y fue validada por el tutor, quien valoró positivamente su enfoque competencial y su conexión con temas de actualidad.

La propuesta encajó perfectamente dentro de los objetivos del módulo, trabajando competencias técnicas, actitudes profesionales y elementos transversales como la seguridad digital. La planificación fue realista en tiempos y recursos, aunque la práctica evidenció la necesidad de prever más profundamente los posibles bloqueos o problemas técnicos derivados de las condiciones del entorno (p.ej., limitaciones del router).

\section{Elaboración y uso de materiales didácticos}

Durante la intervención diseñé y utilicé una presentación propia, con estructura clara, preguntas de reflexión y vídeos breves que facilitaron el debate. También preparé un entorno de red simulado para la parte práctica, lo cual permitió reproducir un contexto de aprendizaje próximo a la realidad profesional.

Los materiales fueron bien acogidos por el alumnado, especialmente por su componente visual y participativo. Además, la nube de palabras inicial y el uso del móvil con fines académicos aportaron un valor añadido a la dinámica.

En futuras intervenciones, sería interesante sistematizar el uso de estos materiales mediante plantillas o rúbricas que permitan replicar y evaluar la experiencia con más precisión.

\section{Valoración del desarrollo de las clases}

Las sesiones se desarrollaron de forma fluida, combinando explicación teórica con actividad práctica. Se fomentó un clima participativo y respetuoso, en el que los estudiantes pudieron expresarse, proponer ideas y colaborar entre sí. Las intervenciones del alumnado fueron en general pertinentes y mostraron interés real por los aspectos técnicos de la seguridad en entornos conectados.

La sesión práctica, en concreto, generó mucho dinamismo y curiosidad. Incluso situaciones no previstas, como el cambio de SSID o el bloqueo del router, se transformaron en oportunidades de aprendizaje valiosas. Esto refuerza la idea de que el aprendizaje en FP debe estar abierto a lo imprevisto, ya que la realidad profesional no siempre sigue un guion fijo.

\section{Evaluación del aprendizaje del alumnado}

Aunque la evaluación formal se limitó a tres preguntas tipo test en el examen del módulo, la actividad permitió observar aprendizajes relevantes no siempre fácilmente evaluables. El alumnado adquirió vocabulario técnico nuevo, comprendió la importancia de la seguridad en sistemas conectados y mostró capacidad de aplicar procedimientos de configuración con criterio.

De cara a futuras intervenciones y diseño de unidades de trabajo, se podría complementar esta evaluación con una rúbrica sencilla o una autoevaluación del alumnado para recoger su percepción del aprendizaje y su implicación.

\section{Consideración sobre criterios DUA}

Durante esta unidad no se aplicaron criterios explícitos del Diseño Universal para el Aprendizaje (DUA). No obstante, sí se favoreció el acceso a la información por múltiples vías (oral, visual, práctica), se permitió el trabajo en grupo y se ofrecieron apoyos individuales durante la actividad. En un futuro, sería interesante incorporar elementos del DUA de forma más sistemática, como permitir distintas formas de expresión o adaptar tareas según niveles de dominio técnico.

\section{Propuestas y sugerencias de mejora}

A partir de esta experiencia, propongo las siguientes mejoras tanto a nivel didáctico como organizativo:

\begin{itemize}
  \item Incluir más contenidos sobre ciberseguridad en el currículo de domótica, no como unidad aislada sino de forma transversal.
  \item Facilitar formación específica al profesorado en metodologías activas como el aprendizaje basado en retos o el uso educativo de tecnologías personales.
  \item Potenciar la reflexión ética y crítica en el uso de dispositivos conectados, incluyendo casos reales, dilemas o debates.
  \item Mejorar la asignación de plazas de prácticas, procurando una correspondencia más ajustada entre el perfil del estudiante y el módulo en el que se le asigna.
  \item Documentar buenas prácticas como esta unidad de trabajo, para que puedan ser replicadas y adaptadas por otros docentes.
\end{itemize}

    % AUTOEVALUACIÓN Y CONCLUSIONES
\chapter{Autoevaluación y Conclusiones}
La experiencia me permitió conocer de cerca la FP y sus metodologías. La principal dificultad fue la desconexión entre mi formación y el contenido impartido. Aun así, logré aportar en el ámbito de la seguridad en IoT y aprender sobre gestión del aula.

\textbf{Conclusión}: El Prácticum II me ha permitido desarrollar competencias pedagógicas clave y reflexionar sobre la importancia de una mejor planificación en la asignación de prácticas.


    % Bibliografía
    \printbibliography
    \nocite{*}


    % Anexos
    \chapter*{Anexos}
    \addcontentsline{toc}{chapter}{Anexos} % Agrega "Anexos" al índice

    \appendix
    \renewcommand{\chaptername}{Anexo} % Cambia "Capítulo" por "Anexo"
    \renewcommand{\thechapter}{\Alph{chapter}} % Usa A, B, C en vez de 1, 2, 3

    \includepdf[
  pages=1,
  scale=0.65,
  offset=0 -3cm,
  pagecommand={
    \chapter{Examen}
    \label{anexo:examen}
  }
]{resources/examen.pdf}

\includepdf[
  pages=2-,
  scale=0.65
]{resources/examen.pdf}

\end{document}










% % Intervención
% \chapter{Intervención y actuación docente en la especialidad}


%     \section{Situación de aprendizaje}
%     Elaborar e implementar una situación de aprendizaje.

%     \begin{table}[]
%         \centering
%         \begin{tabular}{| p{0.6\linewidth} | p{0.35\linewidth} |}
%         \rowcolor[HTML]{86F1E0} 
%         \hline
%         \textbf{Módulo} & \textbf{Curso/Grupo}   \\
%         \hline
%         Instalaciones Domóticas del Ciclo de Técnico de Grado Medio en Instalaciones eléctricas y automáticas (familia Electricidad y Electrónica) & 2º curso / grupo de mañana \\
%         \hline
%         \end{tabular}
%     \end{table}


%     \begin{table}[]
%       \centering
%       \begin{tabular}{| p{0.95\linewidth} |}
%       \rowcolor[HTML]{86F1E0} 
%       \hline
%       \textbf{Título de la unidad de trabajo}   \\
%       \hline
%       Domótica de bajo coste: ¿comodidad inteligente o riesgo silencioso? \\
%       \hline
%       \end{tabular}
%     \end{table}


%     \begin{table}[]
%         \centering
%         \begin{tabular}{| p{0.95\linewidth} |}
%         \rowcolor[HTML]{86F1E0} 
%         \hline
%         \textbf{Contextualización}   \\
%         \hline
%         \textbf{Entorno}: San Vicente del Raspeig es una localidad que forma parte del área metropolitana de Alicante, caracterizada por un entorno urbano a la par que un área industrial con demanda de técnicos electricistas. Además, el auge de la domótica a nivel particular, hace que se incremente la instalación de dispositivos domótcos de bajo coste en viviendas familiares sin tener en cuenta los riesgos de seguridad que puede conllevar este tipo de instalaciones.

%         \textbf{Centro educativo}: El CIPFP Canastell es un centro intregrado de ciclos de formación profesional centrado en el tejido empresarial de alrededores, en concreto en los polígonos industriales.

%         \textbf{Del alumnado y grupo-clase}: el grupo de estudiantes del módulo Instalaciones Domóticas del ciclo Instalaciones eléctricas y automáticas, turno de mañana, está compuesto por alumnado entre 17-25 años aproximadamente. El número de alumnos matriculados es de ¿28?, sin embargo, en el segundo trimestre el nivel de asistencia ha disminuido teniendo un número de presencialidad de unos 12-15 alumnos aprox. Hay alumnos que dan por perdida la asignatura aunque siguen asistiendo a clase. Otros, tienen un nivel mucho más bajo con respecto a los que van más adelantados.

%         \\
%         \hline
%         \end{tabular}
%     \end{table}


%     \begin{table}[]
%         \centering
%         \begin{tabular}{| p{0.95\linewidth} |}
%         \rowcolor[HTML]{86F1E0} 
%         \hline
%         \textbf{Justificación}   \\
%         \hline
%         Esta unidad de trabajo se centra en dar una visión actual de los riesgos de seguridad que pueden tener los dispositivos domóticos de bajo coste, analizando los posibles ataques, vulnerabilidades y cómo reducir éstos aplicando buenas prácticas. De esta forma, el alumnado puede adquirir conocimientos suficientes para asesorar a posibles clientes que les soliciten asesoramiento para realiar instalaciones domóticas.
%         \\
%         \hline
%         \end{tabular}
%     \end{table}


%     \begin{table}[]
%         \centering
%         \begin{tabular}{| p{0.95\linewidth} |}
%         \rowcolor[HTML]{86F1E0} 
%         \hline
%         \textbf{Introducción}   \\
%         \hline
        
%         \\
%         \hline
%         \end{tabular}
%     \end{table}

%     \begin{table}[]
%         \centering
%         \begin{tabular}{| p{0.95\linewidth} |}
%         \rowcolor[HTML]{86F1E0} 
%         \hline
%         \textbf{Objetivos generales}   \\
%         \hline
        
%         \\
%         \hline
%         \end{tabular}
%     \end{table}

%     \begin{table}[]
%         \centering
%         \begin{tabular}{| p{0.95\linewidth} |}
%         \rowcolor[HTML]{86F1E0} 
%         \hline
%         \textbf{Objetivos didácticos}   \\
%         \hline
        
%         \\
%         \hline
%         \end{tabular}
%     \end{table}

%     \begin{table}[]
%         \centering
%         \begin{tabular}{| p{0.95\linewidth} |}
%         \rowcolor[HTML]{86F1E0} 
%         \hline
%         \textbf{Competencias generales}   \\
%         \hline
        
%         \\
%         \hline
%         \end{tabular}
%     \end{table}

%     \begin{table}[]
%         \centering
%         \begin{tabular}{| p{0.95\linewidth} |}
%         \rowcolor[HTML]{86F1E0} 
%         \hline
%         \textbf{Competencias profesionales, personales y sociales}   \\
%         \hline
        
%         \\
%         \hline
%         \end{tabular}
%     \end{table}

%     \begin{table}[]
%         \centering
%         \begin{tabular}{| p{0.95\linewidth} |}
%         \rowcolor[HTML]{86F1E0} 
%         \hline
%         \textbf{Resultados de aprendizaje}   \\
%         \hline
        
%         \\
%         \hline
%         \end{tabular}
%     \end{table}

%     \begin{table}[]
%         \centering
%         \begin{tabular}{| p{0.95\linewidth} |}
%         \rowcolor[HTML]{86F1E0} 
%         \hline
%         \textbf{Criterios de evaluación}   \\
%         \hline
        
%         \\
%         \hline
%         \end{tabular}
%     \end{table}

%     \begin{table}[]
%         \centering
%         \begin{tabular}{| p{0.95\linewidth} |}
%         \rowcolor[HTML]{86F1E0} 
%         \hline
%         \textbf{Contenidos}   \\
%         \hline
        
%         \\
%         \hline
%         \end{tabular}
%     \end{table}

%     \begin{table}[]
%         \centering
%         \begin{tabular}{| p{0.95\linewidth} |}
%         \rowcolor[HTML]{86F1E0} 
%         \hline
%         \textbf{Elementos transversales}   \\
%         \hline
        
%         \\
%         \hline
%         \end{tabular}
%     \end{table}

%     \begin{table}[]
%         \centering
%         \begin{tabular}{| p{0.95\linewidth} |}
%         \rowcolor[HTML]{86F1E0} 
%         \hline
%         \textbf{Metodología}   \\
%         \hline
        
%         \\
%         \hline
%         \end{tabular}
%     \end{table}

%     \begin{table}[]
%         \centering
%         \begin{tabular}{| p{0.95\linewidth} |}
%         \rowcolor[HTML]{86F1E0} 
%         \hline
%         \textbf{Sesión}   \\
%         \hline
        
%         \textbf{Objetivo de la sesión}: \\

%         \textbf{Duración}: \\

%         \textbf{Actividades}:
%         \begin{itemize}
%             \item \textbf{Actividad 1}: Presentación de conceptos, características de los dispositivos domóticos de bajo coste, reflexión sobre vulnerabiidades y consejos básicos para reducir riesgos. Será un trabajo en grupo-clase donde todo el alumnado participará en las reflexiones y cuestiones planteadas. La presentación será proyectada en la pizarra electrónica, teniendo varios puntos de interacción con el alumando para mantener el interés y la atención: al inicio de la presentación, y a modo de rompehielo, se utilizará un código QR a través del cual los alumnos podrán participar en una nube de palabras a través de su dispositvo móvil; en otro momento, se lanzarán interrogantes sobre cuestiones relacionadas con los riesgos de seguridad y saldrán a la pizarra para apuntar lo que se les ha ocurrido.

%             \item \textbf{Actividad 2}: Actividad Práctica de una instalación domótica básica. Con unos pocos dispositivos (como un router, y una bombilla inteligente o similar) se intentará configurar una red simulando una red insegura donde los alumnos podrán "hackear" la red y aprender por qué es importante tener en consideración los aspectos de seguridad. Intentarás conectar la bombilla inteligente a la red, pasando por la instalación de la app correspondiente que les solicitará permisos de WiFi, Bluetooth, Localización, etc. Se realizará con un móvil que no sea del propio alumnado para no comprometer su provacidad. De esta forma, podrán analizar y reflexionar sobre por qué pueden estas aplicaciones solicitar dichos datos, si son seguros, si respetan la privacidad, etc.
%             \begin{itemize}
%                 \item \textbf{Prueba 1}: router huawei (vodafone 2012) no conectado a Internet, configurado con un SSID denominado "RedDomotica" con contraseña débil (12345678). Bombilla inteligente y enchufe inteligente obramart. 
%                 \begin{itemize}
%                     \item Instalación de la app Tuya Smart (se requiere registro)
%                     \item Instalación de la app NavigationSmartHome (para la bombilla) (se requiere registro)
%                     \item La instalación de las app requiere acceso a WiFi, Bluetooth y Localización.
%                     \item El móvil se conecta a la red RedDomotica, la cual no tiene acceso a Internet.
%                     \item Se enciende el dispositivo, primero la bombilla y luego el enchufe por separado. Con el móvil conectado a la red 
%                 \end{itemize}

%                 \item \textbf{Prueba 2}: router huawei (vodafone 2012) configurado como switch, con el mismo SSID y mismo password, ahora con acceso a Internet. Bombilla inteligente y enchufe inteligente obramart. 
%                 \begin{itemize}
%                     \item Instlada la app Tuya Smart
%                     \item El móvil se conecta a la red RedDomotica, la cual no tiene acceso a Internet.
%                     \item Misma operación que en el caso anterior. Parece ser que el indicador de progreso avanza algo más, pero falla igualmente el emparejamiento.
%                 \end{itemize} 

%                 \item \textbf{Prueba 3}:  router tp-link. Configurado con SSID "RedDomotica" y contraseña débil, sin acceso a Internet. Enchufe inteligente Selly Plug y app de la misma marca. 
%                 \begin{itemize}
%                     \item Instalación de la app Shelly (requiere registro: creación de la cuenta clasedomotica6@gmail.com - Qwerty12345!)
%                     \item Se conecta el enchufe inteligente a la corriente. 
%                     \item Se conecta el móvil a RedDomotica
%                     \item Se intenta asignar la red RedSomotica al enchufe.
%                 \end{itemize} 
%             \end{itemize}

%         \end{itemize}
%         \\
%         \hline
%         \end{tabular}
%     \end{table}

% % Autoevaluación y conclusiones
% \chapter{Autoevaluación y Conclusiones}

% % Bibliografía
% \chapter*{Bibliografía}
% \begin{itemize}
%   \item CIPFP Canastell. "Organización del centro". Recuperado de: \url{https://portal.edu.gva.es/cipfpcanastell/informacion-general/organizacion-del-centro/}
%   \item Proyecto de Dirección CIPFP Canastell 2022-2026. Documento interno del centro.
% \end{itemize}



% \end{document}
